\subsection*{怎样把常见消息写短一点}
假设要传递三个符号，出现概率为一半、四分之一、四分之一。统一用两位能表达，但较常出现的符号值得更短的编码。选择 $0,10,11$，读取时不会把某个完整码当成另一码的开头，因而可无歧义地逐个解码。

平均长度是每个码长乘出现概率再求和，本例为一点五位。这里没有任何一个符号使用“一点五位”；分数来自长期平均。这一区别对理解熵与实际整数码长非常重要。

\subsection*{稀有程度为什么变成负对数}
一个公平二选一结果带来一位信息，两个独立结果有四种可能，信息量应相加。概率相乘而信息量相加，提示使用对数；概率越小信息越多，因此取负号。理想长度 $-\log_2p$ 与这种组合规律一致。

熵把各事件的信息量按真实频率平均。二元实验里，接近确定的结果几乎没有意外，公平的一半一半则平均最难预测，所以曲线在中间最高。改变概率前先预测熵的方向，再看动画与数值。

熵给出概率模型下的平均信息尺度，不是某段文字对人的意义，也不是单个字符串必然能压到的长度。单符号前缀码受整数码长约束，长块编码才可逼近更细的平均值；若序列有相关性，还要利用条件分布。下一节问：若编码使用了错误概率，平均要多付多少代价？
